\documentclass[documento]{gpt}
\newcommand{\subsubsubsection}[1]{\paragraph{#1}}
\setcounter{secnumdepth}{4} 
\setcounter{tocdepth}{4}
\usepackage{titlesec}
\titleformat{\paragraph}
  {\normalfont\normalsize\bfseries}
  {\theparagraph}{1em}{}
\titlespacing*{\paragraph}
  {0pt}{3.25ex plus 1ex minus .2ex}{1.5ex plus .2ex}


% ---------------------------------------------------------------------
% Dati specifici del documento
% ---------------------------------------------------------------------
\titolo{Norme di progetto}
\uso{interno}
\destinatari{Prof.~Tullio Vardanega\newline Prof.~Riccardo Cardin}

% Registro delle modifiche: le righe sono nel file
%     norme-di-progetto.registro.tex
% accanto a questo, e le aggiunge il bot dopo il merge in base
% all'etichetta della pull request (GUIDA.md, sezione 5).
% Qui non si scrive nulla.



\begin{document}

% ---------------------------------------------------------------------
% Contenuto del Documento
% ---------------------------------------------------------------------
\section{Introduzione}
\subsection{Scopo del documento}
Questo documento ha lo scopo di definire e descrivere le Norme di Progetto, ovvero l'insieme dei metodi, delle regole e delle procedure operative (il cosiddetto \textit{Way of Working} o \glo{WoW}) adottate dal team per la realizzazione del progetto software. 

L'obiettivo principale della stesura di questo documento è fornire una guida chiara e condivisa per tutti i membri del gruppo. Al suo interno vengono specificati gli strumenti utilizzati e le relative motivazioni, le metodologie di lavoro, gli standard di codifica e le regole di collaborazione, al fine di garantire un approccio strutturato, coerente, efficiente e tracciabile per l'intero ciclo di vita del progetto, limitando il rischio di problemi organizzativi.

La stesura di questo documento segue un approccio incrementale: le norme qui descritte verranno costantemente aggiornate, ampliate e affinate nel tempo per riflettere nuove decisioni tecniche e per supportare il miglioramento continuo dei processi del gruppo.

Per strutturare le attività, il team prende come principale riferimento metodologico le \textit{\glo{best practice}} suggerite dallo \glo{standard} \textbf{ISO/IEC 12207:1995}, adattandole alle specifiche esigenze progettuali. In particolare, lo standard individua tre macro-categorie di processi per il ciclo di vita del software:
\begin{itemize}
    \item \textbf{Processi primari:} includono tutti i processi direttamente coinvolti nella creazione del prodotto software, ovvero le attività di fornitura e sviluppo.
    \item \textbf{Processi di supporto:} comprendono i processi che assistono e sostengono i processi primari garantendone la qualità, la correttezza e la tracciabilità. Tra questi vi sono la documentazione, la gestione della configurazione, l'accertamento della qualità, la verifica e la validazione.
    \item \textbf{Processi organizzativi:} riguardano i processi utilizzati per governare il modo in cui il gruppo lavora, per gestire e coordinare le risorse e per migliorare i flussi operativi. Includono la gestione dei processi, le infrastrutture, la formazione e il miglioramento continuo.
\end{itemize}

\subsection{Glossario}
Al fine di evitare ambiguità e garantire una comprensione chiara e uniforme della terminologia utilizzata, è stato redatto un documento dedicato denominato \textbf{Glossario}. Questo strumento raccoglie e definisce in modo univoco i termini tecnici, gli acronimi, le abbreviazioni e i concetti specifici del dominio del progetto.   Per agevolare la consultazione e la fruizione della documentazione, tutti i vocaboli presenti nel Glossario sono contrassegnati all'interno del testo mediante una ``G'' in pedice (ad esempio: \texttt{termine\textsubscript{G}}). L'impiego del Glossario assicura la coerenza terminologica tra i membri del gruppo e verso gli \glo{stakeholder} esterni, oltre a facilitare l'inserimento operativo di nuovi partecipanti al progetto.  

\section{Processi Primari}

\subsection{Fornitura}
Il processo di \glo{Fornitura} definisce e regolamenta l'insieme delle attività necessarie per avviare, pianificare, realizzare e consegnare il prodotto software in modo controllato, condiviso e verificabile. In qualità di \glo{fornitore}, il gruppo stabilisce le modalità operative, gli strumenti di lavoro e la documentazione contrattuale iniziale per chiarire vincoli, risorse e \glo{requisiti} concordati con la \glo{proponente}. Il processo presidia inoltre il coordinamento delle risorse, il controllo dell'avanzamento e della qualità, nonché la gestione delle comunicazioni verso il \glo{committente} e il supporto nelle fasi di accettazione finale, garantendo la piena conformità del rilascio ai vincoli temporali ed economici pattuiti.

\subsubsection{Documentazione Fornita}
In seguito sono elencati i documenti che, oltre al presente, sono necessari per la consegna del progetto.

\subsubsubsection{Verbali}
I verbali, distinti in interni ed esterni, documentano gli incontri svolti: i primi tengono traccia delle riunioni del team, mentre i secondi riassumono i meeting con l'azienda proponente.

\subsubsubsection{Analisi dei capitolati}
Il documento di Analisi dei Capitolati viene redatto durante la fase di Candidatura con lo scopo di valutare formalmente le proposte d'appalto per il progetto didattico, individuando quella maggiormente allineata alle competenze, agli interessi e agli obiettivi del gruppo.

Per ciascun \glo{capitolato} esaminato, il documento struttura la trattazione secondo i seguenti punti:
\begin{itemize}
    \item \textbf{Descrizione generale:} riassume brevemente il prodotto richiesto e i suoi obiettivi principali;
    \item \textbf{Caratteristiche funzionali e tecnologie:} analizza le funzionalità chiave emerse e lo \glo{stack tecnologico} proposto e/o vincolato;
    \item \textbf{Colloqui informativi:} raccoglie i chiarimenti e i riscontri ottenuti durante gli incontri conoscitivi con l'azienda proponente;
    \item \textbf{Valutazione del gruppo:} esamina i punti di forza, le criticità riscontrate e il livello di gradimento espresso dai membri del team, motivando la scelta finale.
\end{itemize}

\subsubsubsection{Dichiarazione Impegni}
Il documento di Dichiarazione degli Impegni viene redatto durante la fase di Candidatura con lo scopo di formalizzare l'impegno del gruppo nello sviluppo del \glo{capitolato} scelto, definendo la pianificazione economica, oraria e temporale a cui il team si vincola per l'intero \glo{ciclo di vita} del progetto.

In particolare, il documento si articola attraverso i seguenti elementi chiave:
\begin{itemize}
    \item \textbf{Ripartizione dei ruoli e monte ore:} definisce la stima del monte ore complessivo e la sua distribuzione per ciascun ruolo ricoperto dai membri del team;
    \item \textbf{Preventivo dei costi:} determina il costo orario per ruolo e quantifica l'investimento economico complessivo stimato per la realizzazione del prodotto;
    \item \textbf{Scadenza di consegna:} stabilisce la data presunta di rilascio finale del progetto didattico;
    \item \textbf{Formalizzazione dell'impegno:} esplicita la conformità e l'adesione responsabile del gruppo a tali vincoli durante tutte le attività progettuali.
\end{itemize}

\subsubsubsection{Glossario}
Il Glossario costituisce un documento ad aggiornamento continuo lungo l'intero ciclo di vita del progetto, concepito per chiarire e contestualizzare i termini appartenenti al \glo{dominio applicativo} che potrebbero risultare ambigui o di complessa interpretazione. Raccogliendo e formalizzando definizioni per vocaboli tecnici, acronimi e abbreviazioni utilizzati nella documentazione di progetto, esso assicura una comunicazione chiara, coerente e priva di fraintendimenti tra i membri del team, l'azienda proponente, il committente e i lettori esterni.

\subsubsubsection{Lettera di candidatura}
La Lettera di Candidatura rappresenta il documento formale con cui il gruppo GitPushTeam ufficializza la propria candidatura per l'appalto del \glo{capitolato} proposto dall'azienda Eggon. Al suo interno vengono esposte le motivazioni che hanno guidato la decisione del team, valorizzando in particolare l'interesse verso lo stack tecnologico e le opportunità di crescita formativa. Il documento richiama inoltre i principali vincoli operativi assunti dal gruppo tra cui la data di consegna stimata, il preventivo economico e la ripartizione del monte ore, fornendo il necessario raccordo formale con i documenti di Dichiarazione degli Impegni e Analisi dei Capitolati a supporto della proposta.



\section{Processi di Supporto}

\subsection{Documentazione}
Il processo di Documentazione costituisce un processo di supporto fondamentale all'interno del ciclo di vita del software, volto a garantire la tracciabilità, la trasparenza e la continuità di tutte le decisioni prese e delle attività svolte dal gruppo. Essa raccoglie l'insieme delle informazioni scritte necessarie a descrivere il sistema ed i processi, fungendo da mezzo di comunicazione primario sia internamente sia verso l'esterno.

L'obiettivo di questa sezione è definire un insieme di norme, convenzioni e procedure ripetibili atte a:
\begin{itemize}
    \item \textbf{Uniformare la produzione documentale:} stabilire uno stile grafico, strutturale e sintattico omogeneo per tutti gli artefatti del gruppo.
    \item \textbf{Garantire la qualità dei processi:} ridurre ambiguità ed errori tramite regole chiare di stesura e verifica.
    \item \textbf{Semplificare il lavoro del team:} fornire modelli e linee guida operative per agevolare i redattori nella stesura dei documenti.
\end{itemize}

Per la gestione e la produzione della documentazione, il gruppo adotta la filosofia \glo{Documentation as Code}, trattando i sorgenti dei documenti (redatti in LaTeX) alla stesso modo del codice sorgente, avvalendosi dei sistemi di controllo versione  per la loro gestione e tracciabilità. Nelle sottosezioni seguenti vengono nel dettaglio analizzate le regole operative riguardanti la documentazione.

\subsubsection{Template}
Al fine di garantire l'omogeneità grafica, tipografica e strutturale di tutti gli artefatti prodotti dal gruppo, nonché di agevolare le attività di stesura da parte dei redattori, sono stati realizzati diversi modelli personalizzati. Tutti i file di supporto e i sorgenti dei template sono conservati all'interno della directory \texttt{template/modelli} nel repository GitHub del gruppo.

I principali modelli disponibili per la redazione comprendono:

\begin{itemize}
\item \texttt{verbale-interno.tex}: template con la struttura per la documentazione degli incontri interni del gruppo.
\item \texttt{verbale-esterno.tex}: template con la struttura per la documentazione degli incontri esterni tra il gruppo e il proponente.
\item \texttt{lettera.tex}: template utilizzato per la stesura della lettera di presentazione (\autoref{sec:lettera_di_presentazione}) e utilizzabile per documentazione in formato lettera.
\item \texttt{documento.tex}: template di documentazione generica utilizzato nel caso in cui un template più specifico non sia presente.
\end{itemize}

\subsubsection{Versionamento}
Per il versionamento dei nostri documenti e codice abbiamo adottato una versione generalizzata della formula del \glo{versionamento semantico} seguendo lo schema:
\begin{center}
  \Large \texttt{\textbf{Major.Minor.Patch}}
\end{center}

in cui ciascuna cifra risponde al seguente criterio generale:

\begin{itemize}
    \item \textbf{\textit{Major}:} incrementato dalla persona con ruolo di \textit{Responsabile} in sede di approvazione formale.
    \item \textbf{\textit{Minor}:} incrementato in presenza di una modifica sostanziale.
    \item \textbf{\textit{Patch}:} incrementato in caso di modifiche leggere, come la correzione di refusi o typo.
\end{itemize}

Le regole operative che disciplinano il flusso di lavoro all'interno del registro delle modifiche sono le seguenti:
\begin{itemize}
    \item \textbf{Esecuzione delle modifiche:} Chiunque nel gruppo, ad esclusione del \textit{Verificatore}, può effettuare una modifica.
    \item \textbf{Verifica:} Ogni modifica deve essere verificata dal \textit{Verificatore}, il quale, dopo aver completato il controllo, aggiunge il proprio nome al campo verificatore. Questa operazione risulta particolarmente necessaria nel caso in cui in quel momento esistano più verificatori.
    \item \textbf{Approvazione:} Quando il documento raggiunge un punto di maturazione adatto (ad esempio prima di essere inviato al proponente o al committente), la persona che ricopre il ruolo di \textit{Responsabile} approva formalmente il documento e incrementa la versione \textit{Major}.
\end{itemize}

\subsubsection{Nomenclatura}
Al fine di garantire un'identificazione univoca, ordinata ed immediata di ciascun artefatto, il gruppo adotta una convenzione di nomenclatura standardizzata per la denominazione dei file. La struttura del nome varia in base alla tipologia di documento.

\paragraph{Verbali}
I documenti appartenenti alla categoria dei verbali vengono identificati in modo univoco secondo il seguente schema di nomenclatura:

\begin{center}
    \sffamily\bfseries VX-YYYY-MM-DD
\end{center}

in cui:
\begin{itemize}
    \item \textbf{\textit{V}:} rappresenta l'appartenenza alla categoria di documentazione dei verbali.
    \item \textbf{\textit{X}:} rappresenta la categoria di verbale, sarà presente una I in caso di verbale interno oppure una E in caso di verbale esterno.
    \item \textbf{\textit{YYYY-MM-DD}:} indica la data di svolgimento della riunione espressa nel \glo{formato ISO} (\textit{Anno-Mese-Giorno}).
\end{itemize}

\subsection{Gestione della configurazione}
\subsubsection{Identificazione}
\subsubsection{Controllo}
\subsubsection{Registrazione}

% GESTIONE DEI PROCESSI (ALESSIO)
\section{Processi Organizzativi}

\subsection{Gestione dei Processi}
La Gestione dei Processi si prefigge lo scopo di stabilire e standardizzare le metodologie operative ottimali per organizzare il lavoro del team. Attraverso un'attenta supervisione e la regolamentazione dei flussi delle attività, il gruppo punta a prevenire colli di bottiglia e a massimizzare la resa delle risorse a disposizione. Un pilastro di questo processo è l'accurata tracciabilità documentale, necessaria per conservare la memoria delle decisioni e per supportare un'evoluzione qualitativa del prodotto rilasciato.

\subsubsection{Attività previste}
Il processo si articola nelle seguenti macro-attività:
\begin{itemize}
    \item \textbf{Definizione e Inizializzazione:} identificazione del perimetro di lavoro, analisi preliminare degli obiettivi da traguardare e delineazione delle regole base per l'intervento;
    \item \textbf{Programmazione delle attività:} allocazione delle risorse umane, stima delle tempistiche e redazione di un piano strategico per fronteggiare i rischi e monitorare le scadenze;
    \item \textbf{Svolgimento e Monitoraggio:} implementazione pratica dei compiti assegnati, accompagnata da una verifica costante dell'avanzamento lavori per assicurare la massima conformità organizzativa;
    \item \textbf{Valutazione in itinere:} analisi critica dei risultati parziali, volta ad accertare l'aderenza alle specifiche e a integrare eventuali strategie di miglioramento;
    \item \textbf{Conclusione:} formalizzazione del termine dei lavori, aggiornamento delle \glo{baseline} di progetto e consegna del materiale definitivo.
\end{itemize}

\subsubsection{Ruoli di progetto}
Al fine di garantire una formazione trasversale e una distribuzione equa dei carichi, ogni componente del gruppo ricoprirà a rotazione i sei ruoli previsti dal \glo{Way of Working}. Ciascun ruolo è caratterizzato da precise responsabilità e mansioni operative, fondamentali per la buona riuscita del progetto. L'assegnazione avverrà all'inizio di ogni \glo{sprint} e rifletterà le esigenze specifiche di quella particolare fase di sviluppo.

\paragraph{Responsabile}
Questa figura detiene il coordinamento esecutivo del gruppo, assicurandosi che tutte le mansioni convergano verso gli obiettivi prefissati. Stabilisce le scadenze, orchestra le fasi di lavoro e funge da portavoce ufficiale nei confronti della \glo{proponente}. È suo onere, inoltre, anticipare e mitigare eventuali rischi che potrebbero compromettere il progetto.
I suoi compiti sono:
\begin{itemize}
    \item Curare le relazioni e le comunicazioni esterne;
    \item Aggiornare costantemente il \textit{Piano di Progetto};
    \item Redigere i Verbali formali per gli incontri interni ed esterni;
    \item Indire le riunioni, coordinare i lavori e gestire i diari di bordo.
\end{itemize}

\paragraph{Amministratore}
L'Amministratore vigila sull'efficienza e sull'integrità dell'ecosistema di lavoro. Predispone le tecnologie necessarie, configura gli ambienti e si accerta che tutti i membri possano operare senza ostacoli tecnici. Mantiene l'infrastruttura documentale e supporta il \glo{team} nell'adozione degli strumenti scelti.
I suoi compiti sono:
\begin{itemize}
    \item Strutturare e manutenere i \glo{repository} GitHub;
    \item Distribuire il carico di lavoro tramite l'assegnazione delle \glo{issue};
    \item Gestire i canali e le piattaforme di comunicazione interna;
    \item Redigere, aggiornare e perfezionare le \textit{Norme di Progetto} e il \textit{Piano di Qualifica}.
\end{itemize}

\paragraph{Analista}
L'Analista funge da ponte tra le esigenze del committente e il reparto tecnico. Si occupa di recepire, destrutturare e formalizzare le richieste, traducendole in specifiche chiare, inequivocabili e pronte per essere elaborate nelle successive fasi di \glo{design}.
I suoi compiti sono:
\begin{itemize}
    \item Approfondire la conoscenza del dominio applicativo e dei relativi vincoli di business;
    \item Interfacciarsi con la \glo{proponente} per validare i requisiti estratti;
    \item Curare la stesura del documento di \textit{Analisi dei Requisiti}.
\end{itemize}

\paragraph{Progettista}
Questa figura ha il compito di delineare l'architettura logica e tecnologica del prodotto. Partendo dall'output degli analisti, il Progettista seleziona le tecnologie idonee e modella la struttura del \glo{software}, garantendo scalabilità, robustezza e modularità al sistema.
I suoi compiti sono:
\begin{itemize}
    \item Valutare e selezionare lo \glo{stack tecnologico} ottimale, pesandone pro e contro;
    \item Definire l'impianto architetturale e le logiche di interazione tra i vari componenti;
    \item Vigilare affinché l'implementazione pratica del codice rispetti i \glo{design pattern} e l'architettura delineata.
\end{itemize}

\paragraph{Programmatore}
Il Programmatore traduce le direttive del Progettista in codice eseguibile. Si concentra sull'implementazione pratica delle funzionalità, mantenendo un'elevata attenzione alla qualità del codice sorgente, alla sua manutenibilità e all'aderenza agli standard di progetto.
I suoi compiti sono:
\begin{itemize}
    \item Sviluppare il codice sorgente nel rigoroso rispetto delle \textit{Norme di Progetto};
    \item Individuare e correggere prontamente i \glo{bug} emersi in fase di esecuzione;
    \item Realizzare adeguate suite di \glo{test} per validare il comportamento e la sicurezza dei moduli sviluppati.
\end{itemize}

\paragraph{Verificatore}
Il Verificatore agisce come garante ultimo della qualità del progetto. Sottopone a rigorosi controlli ogni \glo{artefatto} prodotto dal team, che si tratti di documentazione o codice sorgente, accertandosene della conformità prima dell'approvazione e del rilascio finale.
I suoi compiti sono:
\begin{itemize}
    \item Ispezionare i documenti redatti per assicurarne coerenza strutturale, assenza di refusi e correttezza formale;
    \item Eseguire procedure di verifica strutturale e \glo{code review} per garantire un \glo{software} affidabile, ben commentato e stilisticamente coerente.
\end{itemize}
\clearpage
\subsection{Comunicazione}
Verranno di seguito elencate le modalità di comunicazione, sincrone e non, che verranno adottate dal gruppo per mantenere un contatto proficuo e attivo tra i propri componenti interni e i soggetti esterni.
\paragraph{Comunicazioni interne}
Per quanto riguarda la comunicazione interna, il gruppo ha deciso di adottare le seguenti modalità:
\begin{itemize}
\item Asincrone: \begin{itemize}
    \item \textbf{WhatsApp}: è stato creato un gruppo WhatsApp per il progetto, ove tutti i membri del gruppo sono presenti, ed è stato scelto, in accordo con tutti i componenti del team, di renderlo il canale principale per comunicazioni asincrone e organizzative. La motivazione della scelta è data dalla versatilità dello strumento, la familiarità che tutti i membri del gruppo hanno in esso, e la sua estrema semplicità d'uso. Le attività principali che vengono svolte in questo canale sono:\begin{itemize}
        \item Discussioni semplici affrontabili in modalità differita.
        \item Organizzazione di incontri straordinari o modulazione degli orari standard.
        \item Chiarimento di dubbi e difficoltà progettuali che non richiedono un incontro sincrono.
        \item Recap delle principali attività assegnate ad ogni membro del team dopo una riunione interna.
        \end{itemize}
    \item \textbf{Chat testuale di Discord}: utilizzata per scambio di informazioni utili allo svolgimento di determinate task (nel canale generale) e per scambio di informazioni e documenti durante fasi di lavoro collaborativo (nei sotto-canali dedicati).
        \end{itemize}
    \item Sincrone: \begin{itemize}
        \item \textbf{Discord}: è la piattaforma che è stata scelta come mezzo per svolgere le riunioni interne. Discord consente di effettuare sia chiamate vocali standard che videoconferenze, con anche la possibilità di condividere lo schermo durante le fasi di presentazione. Il canale inoltre è stato suddiviso in sotto-canali, dove parte dei membri del team possono trovarsi per sessioni di lavoro collaborativo e discussioni interne, sono inoltre presenti dei canali testuali, già descritti nella parte precdente.
        \item \textbf{Incontri in presenza}: Si è scelto di non effettuare incontri in presenza con cadenza standard, ma solo in situazioni straordinarie decise in comunione dall'intero team.
    \end{itemize}
    \end{itemize}
\paragraph{Riunioni Interne}
Come previo descritto, le riunioni interne verranno svolte in remoto via Discord, preferito per la facilità di organizzazione e la maggiore integrazione con gli impegni e le necessità dei singoli membri del team, anche per la distanza geografica presente tra i singoli membri.
Verranno svolte:\begin{itemize}
    \item Una riunione \textbf{interna} dalla cadenza settimanale il lunedì, con scelta dell'orario variabile tra il tardo pomeriggio e la sera, discussa di volta in volta mediante i canali di comunicazione asincroni.
    \item Riunioni \textbf{straordinarie} qualora si ritenesse necessario un ulteriore incontro causato da difficoltà critiche o necessità di lavoro collaborativo per moli di lavoro straordinarie.
\end{itemize}
Nel caso delle riunioni straordinarie, la data e l'ora verrà discussa su WhatsApp ed adattata in base alle esigenze dei membri. Le riunioni settimanali non vedranno rimodulazione di data se non in casi straordinari dati da impegni improrogabili di due o più membri del team.
\paragraph{Comunicazioni Esterne}
Da definire col proponente una volta assegnato l'appalto.
\paragraph{Riunioni Esterne con Proponenete e Committente}
Da definire.
\paragraph{Pianificazione degli Impegni}
Per semplificare la pianificazione degli impegni dei singoli membri del team, per incontri straordinari e non, si utilizzaono principalmente dei sondaggi su WhatsApp di volta in volta per definire orari consoni con gli impegni dei singoli membri. In fase di valutazione l'utilizzo di un calendario condiviso (Es. Google Calendar dell'account del team) ove esplicitare formalmente gli impegni settimanali dei vari membri.

\subsection{Infrastruttura}
Il processo di Infrastruttura ha lo scopo di predisporre e mantenere nel tempo
l'insieme di strumenti, servizi e ambienti di lavoro di cui hanno bisogno tutti
gli altri processi: dalla stesura della documentazione alla sua verifica e
pubblicazione, dalla pianificazione delle attività alla comunicazione interna
ed esterna. Un'infrastruttura adeguata permette a ogni membro del gruppo di
lavorare senza ostacoli tecnici, con gli stessi strumenti e le stesse regole,
qualunque sia il ruolo che ricopre in quel momento.
 
La responsabilità del processo è dell'\textbf{Amministratore}, che crea,
configura e mantiene gli strumenti. Ogni membro del gruppo è tenuto a usarli
secondo quanto descritto in queste Norme e a segnalare tempestivamente
qualsiasi malfunzionamento.
 
\subsubsection{Attività previste}
Le attività previste dal processo di Infrastruttura sono:
\begin{itemize}
  \item \textbf{Implementazione del processo:} individuazione dei bisogni dei
        processi da supportare, scelta degli strumenti che li soddisfano e
        definizione delle regole con cui devono essere usati. Le scelte sono
        documentate nel presente documento;
  \item \textbf{Creazione dell'infrastruttura:} configurazione degli
        strumenti scelti e loro messa a disposizione di tutto il gruppo,
        verificandone il funzionamento prima dell'adozione;
  \item \textbf{Manutenzione dell'infrastruttura:} aggiornamento,
        correzione ed evoluzione degli strumenti, così che restino
        disponibili e adeguati alle esigenze del progetto in ogni sua fase.
\end{itemize}
 
\subsubsection{Principi adottati}
Nella scelta e nella configurazione degli strumenti il gruppo segue questi
principi:
\begin{itemize}
  \item \textbf{Un'unica fonte ufficiale:} il \glo{repository} della
        documentazione è l'unico luogo in cui risiede la versione ufficiale di
        ogni documento. Bozze e appunti possono nascere altrove (si veda
        \S~\ref{sec:strumenti-collaborazione}), ma un contenuto diventa
        ufficiale solo quando entra nel repository tramite \glo{pull request}
        verificata;
  \item \textbf{Automazione dei controlli ripetitivi:} compilazione,
        controlli di forma e pubblicazione sono svolti in automatico, così
        da ridurre gli errori manuali e lasciare alla verifica umana il
        controllo dei contenuti;
  \item \textbf{Nessuna modifica diretta al ramo principale:} il ramo
        \texttt{main} è protetto e si aggiorna solo tramite pull request
        approvata da un membro diverso dall'autore;
  \item \textbf{Strumenti gratuiti e accessibili:} sono adottati solo
        strumenti senza costi, disponibili su tutti i sistemi operativi usati
        dai membri del gruppo;
  \item \textbf{Riproducibilità:} un documento produce lo stesso risultato
        sul computer di ciascun membro e nel sistema di integrazione
        continua, perché template, configurazione e script di compilazione
        sono versionati insieme ai documenti.
\end{itemize}
 
\subsubsection{Infrastruttura del gruppo}
L'infrastruttura del gruppo è costituita dai seguenti elementi.
 
\paragraph{Organizzazione GitHub}
Tutte le risorse del progetto risiedono nell'organizzazione GitHub
\textbf{Git-Push-Team} (\url{https://github.com/Git-Push-Team}), di cui fanno
parte tutti i membri del gruppo.
 
\paragraph{Repository della documentazione}
Il repository \texttt{docs} (\url{https://github.com/Git-Push-Team/docs})
contiene i sorgenti \LaTeX{} di tutti i documenti, organizzati per milestone:
\begin{itemize}
  \item \texttt{Candidatura/}, \texttt{RTB/} e \texttt{PB/}: una cartella per
        milestone, ciascuna suddivisa in \texttt{documenti/} e
        \texttt{verbali/}, a loro volta divise in \texttt{interni/} ed
        \texttt{esterni/};
  \item \texttt{template/}: la classe \LaTeX{} comune \texttt{gpt.cls}, i dati
        del gruppo, i modelli da cui partire per ogni nuovo documento e alcuni
        esempi completi;
  \item \texttt{scripts/}: gli script Python che compilano i documenti, ne
        controllano la forma e generano il sito;
  \item \texttt{.github/}: le automazioni, i modelli delle \glo{issue} e il
        modello delle pull request;
  \item \texttt{GUIDA.md}: il manuale operativo che descrive passo per passo
        come svolgere ogni attività sul repository.
\end{itemize}
 
\paragraph{Integrazione continua}
Il repository usa GitHub Actions per tre automazioni:
\begin{itemize}
  \item \textbf{Anteprima delle pull request:} a ogni modifica proposta
        vengono compilati tutti i documenti e nella pull request compare un
        commento con il PDF dei documenti modificati, gli eventuali errori di
        compilazione e gli avvisi sui controlli automatici (registro delle
        modifiche, termini del Glossario, posizione dei file);
  \item \textbf{Pubblicazione:} dopo l'unione in \texttt{main} i documenti
        vengono ricompilati, i PDF vengono salvati nel repository accanto ai
        sorgenti e il sito della documentazione viene aggiornato;
  \item \textbf{Rilascio delle milestone:} alla creazione di un'etichetta
        di versione (\textit{tag}) viene prodotta una \textit{release} che
        contiene i PDF della milestone consegnata.
\end{itemize}
 
\paragraph{Sito della documentazione}
I documenti approvati sono pubblicati con GitHub Pages all'indirizzo
\url{https://git-push-team.github.io/docs/}. Il sito espone per ogni
milestone la Lettera di Presentazione, i documenti e i verbali, sempre
nell'ultima versione, e la pagina del Glossario. Ogni risorsa si apre in una
scheda dedicata.
 
\paragraph{Gestione delle attività}
Le attività sono tracciate con le issue del repository e raccolte nel
GitHub Project \textbf{GIT PUSH TEAM}, dove per ognuna sono registrati stato,
priorità, ruolo, ore stimate, ore effettive, data di inizio e scadenza.
Le issue si aprono con due modelli: \textit{Attività}, per un compito da
svolgere su un documento, e \textit{Errore}, per un errore trovato in un
documento.
 
\subsubsection{Strumenti}
Di seguito sono elencati gli strumenti adottati, raggruppati per funzione, con
l'uso che ne fa il gruppo.
 
\paragraph{Gestione della configurazione e del progetto}
\begin{description}
  \item[Git] (\url{https://git-scm.com/}) \hfill \\
        Sistema di controllo di versione distribuito. Ogni attività si svolge
        in un \glo{branch} dedicato, creato a partire dalla issue
        corrispondente.
  \item[GitHub] (\url{https://github.com/}) \hfill \\
        Servizio che ospita i repository del gruppo. Oltre al codice e ai
        documenti, il gruppo ne usa le \textit{pull request} per la verifica,
        le \textit{issue} e i \textit{Projects} per la pianificazione,
        \textit{Actions} per l'integrazione continua e \textit{Pages} per il
        sito.
  \item[GitHub Desktop] (\url{https://desktop.github.com/}) \hfill \\
        Applicazione grafica per usare Git senza riga di comando: permette di
        creare branch, registrare le modifiche e pubblicarle con pochi click.
\end{description}
 
\paragraph{Documentazione}
\begin{description}
  \item[\LaTeX{} e TeX Live] (\url{https://www.tug.org/texlive/}) \hfill \\
        Linguaggio e distribuzione usati per scrivere tutti i documenti. La
        compilazione è affidata a \texttt{latexmk}, sia in locale sia
        nell'integrazione continua.
  \item[Template del gruppo] \hfill \\
        La classe \texttt{gpt.cls} genera da sola copertina, registro delle
        modifiche, indici, intestazioni e piè di pagina, e controlla durante
        la compilazione la correttezza formale di dati come versioni, date e
        nomi dei membri.
  \item[Visual Studio Code] (\url{https://code.visualstudio.com/}) \hfill \\
        Editor consigliato per la stesura, con l'estensione \textit{LaTeX
        Workshop} che compila il documento e ne mostra il PDF a ogni
        salvataggio. Il repository ne contiene già la configurazione.
  \item[Overleaf] (\url{https://www.overleaf.com/}) \hfill \\
        Editor \LaTeX{} online, usato quando più membri devono scrivere
        contemporaneamente sullo stesso testo. Il testo prodotto va poi
        riportato nel repository con una pull request: solo allora diventa
        parte del documento ufficiale.
  \item[Python] (\url{https://www.python.org/}) \hfill \\
        Linguaggio degli script che compilano i documenti, eseguono i
        controlli automatici e generano il sito.
\end{description}
 
\paragraph{Comunicazione}
\begin{description}
  \item[Discord] (\url{https://discord.com/}) \hfill \\
        Piattaforma per le riunioni interne e le sessioni di lavoro
        collaborativo, con canali vocali, testuali e condivisione dello
        schermo.
  \item[WhatsApp] (\url{https://www.whatsapp.com/}) \hfill \\
        Canale principale per le comunicazioni asincrone e organizzative.
  \item[Gmail] (\url{https://mail.google.com/}) \hfill \\
        Ospita l'indirizzo del gruppo, \texttt{gitpushteam@gmail.com}, usato
        per tutte le comunicazioni ufficiali con il committente e con le
        aziende proponenti.
  \item[Google Meet e Microsoft Teams] (\url{https://meet.google.com/},
        \url{https://www.microsoft.com/microsoft-teams/}) \hfill \\
        Piattaforme per le videochiamate con le aziende proponenti. La
        piattaforma è scelta dall'azienda che organizza l'incontro.
\end{description}
 
\paragraph{Collaborazione}\label{sec:strumenti-collaborazione}
\begin{description}
  \item[Google Drive, Documenti e Fogli] (\url{https://drive.google.com/})
        \hfill \\
        Spazio condiviso per bozze, appunti presi durante le riunioni e
        calcoli preliminari, come le stime di ore e costi. Questi file sono
        materiale di lavoro e non sostituiscono i documenti ufficiali: il
        loro contenuto definitivo va riportato nel repository.
\end{description}
 
\subsubsection{Manutenzione e modifica dell'infrastruttura}
\begin{itemize}
  \item Chi riscontra un malfunzionamento di uno strumento, del template o
        delle automazioni lo segnala con una issue, descrivendo il problema e
        come riprodurlo.
  \item Le modifiche a template, script e automazioni sono svolte
        dall'Amministratore in un branch dedicato e passano, come ogni altra
        modifica, da una pull request verificata da un altro membro. Prima
        di essere unite, quelle di maggiore impatto sono provate in un
        repository di prova, per non introdurre errori nel repository
        ufficiale.
  \item Ogni modifica che cambia il modo di lavorare del gruppo viene
        descritta nel manuale operativo \texttt{GUIDA.md} e, se riguarda le
        regole, anche nelle presenti Norme. Il resto del gruppo viene
        informato nella riunione successiva.
  \item L'adozione di un nuovo strumento, o l'abbandono di uno esistente, è
        una decisione del gruppo: viene presa in riunione e registrata nel
        verbale.
\end{itemize}
 
% ---------------------------------------------------------------------
 
\subsection{Miglioramento}
Il processo di Miglioramento ha lo scopo di valutare in modo sistematico i
processi adottati dal gruppo e di modificarli quando non risultano efficaci o
efficienti. Le Norme di Progetto non sono quindi un documento definitivo: si
aggiornano ogni volta che il gruppo individua un modo migliore di lavorare, e
ogni aggiornamento è tracciato nel registro delle modifiche.
 
Il processo segue il ciclo \textbf{PDCA} (\textit{Plan, Do, Check, Act}), che
viene ripetuto a ogni sprint:
\begin{itemize}
  \item \textbf{Plan:} si individua un problema o un'opportunità e si
        pianifica l'azione per affrontarlo;
  \item \textbf{Do:} si attua l'azione pianificata;
  \item \textbf{Check:} si verifica se l'azione ha prodotto l'effetto atteso;
  \item \textbf{Act:} se l'esito è positivo, la nuova pratica diventa una
        regola e si registra nelle Norme; altrimenti si corregge l'azione e
        il ciclo riparte.
\end{itemize}
 
\subsubsection{Attività previste}
Le attività previste dal processo di Miglioramento sono:
\begin{itemize}
  \item \textbf{Implementazione del processo:} definizione dei processi del
        gruppo, delle occasioni in cui vengono valutati e di chi se ne
        occupa;
  \item \textbf{Valutazione del processo:} raccolta di dati e riscontri
        sull'andamento dei processi, per individuarne punti di forza e
        criticità;
  \item \textbf{Miglioramento del processo:} definizione, attuazione e
        controllo delle azioni correttive, fino al loro consolidamento nelle
        Norme.
\end{itemize}
 
\subsubsection{Implementazione del processo}
I processi del gruppo sono formalizzati nel presente documento, che ne
costituisce il riferimento, e nel manuale operativo \texttt{GUIDA.md} per gli
aspetti pratici. Essendo versionati nel repository, ogni loro modifica è
tracciabile: si sa che cosa è cambiato, quando, chi l'ha proposta e chi l'ha
verificata.
 
I ruoli coinvolti nel processo sono:
\begin{itemize}
  \item il \textbf{Responsabile}, che conduce la retrospettiva, raccoglie le
        proposte di miglioramento e ne pianifica l'attuazione;
  \item l'\textbf{Amministratore}, che attua le modifiche a strumenti e
        procedure e aggiorna le Norme di Progetto e la \texttt{GUIDA.md};
  \item il \textbf{Verificatore}, che controlla che le modifiche siano
        coerenti con il resto delle Norme;
  \item \textbf{ogni membro del gruppo}, che può proporre un miglioramento in
        qualsiasi momento, aprendo una issue o portandolo alla riunione
        successiva.
\end{itemize}
 
\subsubsection{Valutazione del processo}
Il gruppo lavora per \glo{sprint} di due settimane, che coincidono con il
periodo di rotazione dei ruoli. La valutazione dei processi avviene nella
\textbf{retrospettiva}, un punto fisso della riunione del lunedì che chiude
ogni sprint, prima del cambio dei ruoli. Durante la retrospettiva ciascun
membro risponde a tre domande:
\begin{itemize}
  \item \textbf{che cosa ha funzionato}, e va quindi mantenuto;
  \item \textbf{che cosa non ha funzionato}, e perché;
  \item \textbf{che cosa cambiamo} nello sprint successivo.
\end{itemize}
 
Per rendere la valutazione oggettiva, la retrospettiva si basa anche su dati
raccolti durante lo sprint:
\begin{itemize}
  \item il confronto tra ore stimate e ore effettive delle attività,
        registrate nel GitHub Project;
  \item le attività concluse in ritardo o rimaste aperte rispetto a quelle
        pianificate;
  \item il numero di correzioni richieste in verifica e gli avvisi ricorrenti
        dei controlli automatici;
  \item i valori delle metriche di processo definite nel Piano di Qualifica,
        quando disponibili.
\end{itemize}
 
Oltre alla retrospettiva, sono occasioni di valutazione i riscontri esterni:
le osservazioni del committente (ad esempio in sede di aggiudicazione e alle
revisioni di avanzamento) e quelle dell'azienda proponente negli incontri
periodici. Ogni osservazione ricevuta viene esaminata nella prima riunione
utile.
 
\subsubsection{Miglioramento del processo}
Ogni criticità emersa dalla valutazione segue lo stesso percorso, così che
nessun miglioramento resti solo un'intenzione:
\begin{enumerate}
  \item la criticità e la soluzione scelta sono registrate come
        \textbf{decisione} nel verbale della riunione;
  \item dalla decisione nasce un'\textbf{azione tracciabile}: una issue con
        un incaricato e una scadenza, riportata anche nel verbale;
  \item l'incaricato attua l'azione e, se cambia una regola, aggiorna le
        Norme di Progetto con una pull request verificata;
  \item nella retrospettiva successiva il gruppo controlla se l'azione ha
        risolto la criticità. In caso positivo la nuova pratica resta in
        vigore; in caso contrario si pianifica un'azione correttiva diversa.
\end{enumerate}
 
\subsubsection{Miglioramenti attuati}
La Tabella~\ref{tab:miglioramenti} raccoglie i miglioramenti già attuati dal
gruppo, con la criticità da cui sono nati e l'azione intrapresa. La tabella
viene aggiornata a ogni nuovo miglioramento.
 
\begin{table}[H]
  \centering
  \small
  \caption{Miglioramenti attuati\label{tab:miglioramenti}}
  \begin{tabularx}{\textwidth}{@{}p{1.9cm}XX@{}}
    \toprule
    \intestazione{Data} & \intestazione{Criticità rilevata} &
      \intestazione{Azione intrapresa} \\
    \midrule
    2026-10-20 &
      Nella prima struttura del repository i documenti erano divisi per
      tipo, mentre la suddivisione per milestone compariva solo sul sito:
      chi cercava un file non lo trovava nello stesso posto nei due luoghi. &
      Riorganizzazione del repository per milestone, poi per documenti e
      verbali, interni ed esterni, come sul sito. La nuova struttura è stata
      provata in un repository di prova prima dell'adozione (verbale
      interno del 2026-10-20). \\
    \addlinespace
    2026-10-29 &
      Dall'esame delle osservazioni del committente ai gruppi degli anni
      precedenti è emerso che alcune righe dei nostri registri delle
      modifiche avevano una versione senza verifica, righe di sola
      ``Verifica'' senza autore e date non coerenti. &
      Nuove regole di versionamento: ogni riga ha autore e verificatore
      distinti, la versione cambia solo dopo una verifica superata e la
      versione 1.0.0 indica l'approvazione. Correzione dei registri esistenti
      e controllo automatico del registro in ogni pull request. \\
    \addlinespace
    2026-10-29 &
      Dalle stesse osservazioni: le risorse di un sito di documentazione
      devono aprirsi in una scheda dedicata, che ne riporti il nome reale. &
      Sul sito ogni documento si apre in una nuova scheda, intitolata con il
      titolo e la versione del documento. \\
    \bottomrule
  \end{tabularx}
  
\end{table}

\end{document}
